\section{\emph{Traditionis Custodes} and the Absence of a Published Special Act}

Francis's \emph{Traditionis custodes}, issued 16 July 2021, replaced the general categories of 2007. It called the liturgical books promulgated by Paul VI and John Paul II the unique expression of the Roman Rite's \emph{lex orandi}, made the diocesan bishop responsible for authorizing use of the 1962 Missal under Apostolic See guidelines, restricted locations and new groups, and abrogated contrary prior norms and customs. Article 6 placed institutes and societies erected by PCED under the competence of the consecrated-life dicastery.

The accompanying papal letter explained a concern for unity and directed bishops toward the reformed liturgy over time. The Divine Worship dicastery's December 2021 responses and a papal rescript publicized in February 2023 further specified reserved decisions in the general regime. These are controlling public acts at their dates. They do not name the ICKSP or print the complete relation between its proper law and every new restriction.

\subsection{Existence, patrimony, and permission are different questions}

The Institute remained a pontifical-right society of apostolic life. \emph{Traditionis custodes} did not excardinate its priests, dissolve its government, or suppress its houses as such. Continued existence, however, is not identical with a universal right to celebrate every public liturgy in every place. Canon 738 and the act's episcopal provisions make local authority materially relevant.

The Institute's 2008 decree had approved a canonial society whose worship was described under the 2007 framework. Definitive constitutions were reported in 2016. How those acts interact with later universal legislation is a legal question that cannot be solved by simply declaring the older act vanished or the newer act irrelevant. Public practice after 2021 shows variation rather than a single transparent rule applied identically everywhere.

\subsection{No checked universal ICKSP exemption}

The FSSP received a papal decree in February 2022 published by that society, treated in TPI-02. No equivalent ICKSP-named universal faculty, exemption, rescript, or later papal act was located in the checked Vatican and institutional sources through 16 July 2026. The absence of a published act must not be converted into proof that no private communication or local permission exists. It does mean this study cannot transfer the FSSP decree to the ICKSP by analogy.

The 2024 ICKSP report of a private papal audience says Francis encouraged service according to the Institute's proper charism. That is important participant evidence of relationship. It is not legislation, a transcript, or an exemption from \emph{Traditionis custodes}. Pastoral encouragement and juridic faculty are different source types.

\claimcheck{``All former-PCED institutes received the FSSP's 2022 exemption.''}{The published FSSP decree names and grants faculties to FSSP members. No checked public instrument extends its operative terms to the ICKSP. Shared historical supervision does not make one society's special act the other's law.}
